\documentclass[12pt]{article}
\usepackage{amsmath}
\usepackage[margin=1in]{geometry}
\begin{document}

\section*{Problem \#1}
Determine the gravitational potential energy, in kJ, of \(2 \, \text{m}^3\) of liquid water at an elevation of \(30 \, \text{m}\) above the surface of Earth. Also, determine the change in potential energy if the elevation decreases by \(15 \, \text{m}\).

\section*{Solution}

\subsection*{Calculating Initial Gravitational Potential Energy}
The formula for gravitational potential energy \(U\) is:

\[
U = m \times g \times h
\]

Here, \(m\) is mass, \(g\) is the acceleration due to gravity (\(9.7 \, \text{m/s}^2\)), and \(h\) is the height above the surface of the Earth (\(30 \, \text{m}\)).

First, we need to find the mass of \(2 \, \text{m}^3\) of liquid water. Since the density \(\rho\) of water is \(1000 \, \text{kg/m}^3\), the mass \(m\) can be calculated as:

\[
m = \rho \times V = 1000 \, \text{kg/m}^3 \times 2 \, \text{m}^3 = 2000 \, \text{kg}
\]

Now, we can find the initial gravitational potential energy \(U_{\text{initial}}\):

\[
U_{\text{initial}} = 2000 \, \text{kg} \times 9.7 \, \text{m/s}^2 \times 30 \, \text{m} = 582000 \, \text{J} = 582 \, \text{kJ}
\]

\subsection*{Calculating the Change in Potential Energy}

To find the change in potential energy, we first find the potential energy at the new elevation \(h' = h - 15\), which is \(30 \, \text{m} - 15 \, \text{m} = 15 \, \text{m}\).

\[
U_{\text{final}} = 2000 \, \text{kg} \times 9.7 \, \text{m/s}^2 \times 15 \, \text{m} = 291000 \, \text{J} = 291 \, \text{kJ}
\]

The change in potential energy \(\Delta U\) is:

\[
\Delta U = U_{\text{final}} - U_{\text{initial}} = 291 \, \text{kJ} - 582 \, \text{kJ} = -291 \, \text{kJ}
\]

Thus, the initial gravitational potential energy is \(582 \, \text{kJ}\) and the change in potential energy due to the decrease in elevation is \(-291 \, \text{kJ}\).


\newpage
\section*{Problem \#2}

Given a vertical piston-cylinder assembly containing air, you are tasked with finding the change in specific internal energy (\( \Delta u \)) under the following conditions:

\begin{itemize}
    \item Mass of the piston: \( m_{\text{piston}} = 25 \, \text{kg} \)
    \item Face area of the piston: \( A = 0.005 \, \text{m}^2 \)
    \item Mass of air: \( m_{\text{air}} = 2.5 \, \text{g} \)
    \item Initial volume of air: \( V_{\text{initial}} = 2.5 \, \text{liters} \)
    \item Atmospheric pressure: \( P_{\text{atm}} = 100 \, \text{kPa} \)
    \item Final volume of air: \( V_{\text{final}} = 0.001 \, \text{m}^3 \)
    \item Energy removed by heat transfer: \( Q = -1 \, \text{kJ} \)
\end{itemize}

We are to neglect friction, potential energy (PE), and kinetic energy (KE).

\section*{Solution}

\subsection*{Calculating Work Done on the System}

First, we need to calculate the work done on the system when the volume changes. The work (\( W \)) done in a quasi-static isochoric process is given by:

\[
W = - \int_{V_{\text{initial}}}^{V_{\text{final}}} P \, dV
\]

Here, \( P \) is the pressure exerted by the piston and the atmosphere on the air inside the cylinder, which can be calculated as:

\[
P = P_{\text{atm}} + \frac{m_{\text{piston}} \cdot g}{A}
\]
  
Substituting \( P_{\text{atm}} = 100 \, \text{kPa} \), \( m_{\text{piston}} = 25 \, \text{kg} \), \( g = 9.81 \, \text{m/s}^2 \), and \( A = 0.005 \, \text{m}^2 \), we get:

\[
P = 100 \, \text{kPa} + \frac{25 \, \text{kg} \times 9.81 \, \text{m/s}^2}{0.005 \, \text{m}^2}
\]
\[
P = 100 \, \text{kPa} + 49.125 \, \text{kPa}
\]
\[
P = 149.125 \, \text{kPa}
\]

Now, we can integrate to find the work done:

\[
W = -149.125 \times 10^3 \, \text{N/m}^2 \times (0.001 \, \text{m}^3 - 0.0025 \, \text{m}^3)
\]
\[
W = 149.125 \, \text{N/m}^2 \times 0.0015 \, \text{m}^3
\]
\[
W = 0.2236875 \, \text{kJ}
\]

\subsection*{Calculating the Change in Internal Energy}

The First Law of Thermodynamics states that \( \Delta U = Q - W \). We already have \( W = 0.2236875 \, \text{kJ} \) and \( Q = -1 \, \text{kJ} \), so:

\[
\Delta U = -1 \, \text{kJ} - 0.2236875 \, \text{kJ}
\]
\[
\Delta U = -1.2236875 \, \text{kJ}
\]

Since \( \Delta U \) is for 2.5 g of air, we need to find the change in specific internal energy (\( \Delta u \)):

\[
\Delta u = \frac{\Delta U}{m_{\text{air}}}
\]

First, convert the mass of air to kg: \( m_{\text{air}} = 2.5 \, \text{g} = 0.0025 \, \text{kg} \)

\[
\Delta u = \frac{-1.2236875 \, \text{kJ}}{0.0025 \, \text{kg}}
\]
\[
\Delta u = -489.475 \, \text{kJ/kg}
\]

Therefore, the change in specific internal energy is \( -489.475 \, \text{kJ/kg} \).
\newpage
\section*{Problem \#3}

A gas in a piston--cylinder assembly undergoes a process for which the relationship between pressure and volume is \( PV^2 = \text{constant} \). The initial state is \( P_{\text{initial}} = 1 \, \text{bar} \) and \( V_{\text{initial}} = 0.1 \, \text{m}^3 \), and the final state is \( P_{\text{final}} = 9 \, \text{bar} \). We are tasked to find:

\begin{enumerate}
    \item The final volume \( V_{\text{final}} \) in \( \text{m}^3 \)
    \item The work done during this process
\end{enumerate}

\section*{Solution}

\subsection*{Calculating the Final Volume}

The relationship between pressure and volume is given as \( PV^2 = \text{constant} \). Thus, we can write:

\[
P_{\text{initial}} \times V_{\text{initial}}^2 = P_{\text{final}} \times V_{\text{final}}^2
\]

Given \( P_{\text{initial}} = 1 \, \text{bar} \), \( V_{\text{initial}} = 0.1 \, \text{m}^3 \), and \( P_{\text{final}} = 9 \, \text{bar} \), we can solve for \( V_{\text{final}} \):

\[
V_{\text{final}} = \sqrt{\frac{P_{\text{initial}} \times V_{\text{initial}}^2}{P_{\text{final}}}} = \sqrt{\frac{1 \, \text{bar} \times (0.1 \, \text{m}^3)^2}{9 \, \text{bar}}} \approx 0.0333 \, \text{m}^3
\]
\subsection*{Calculating the Work Done}

The work \( W \) done during the process is an integral involving pressure \( P \) and a change in volume \( dV \):

\[
W = \int_{V_{\text{initial}}}^{V_{\text{final}}} P \, dV
\]
In this specific process, the relationship between pressure and volume is \( PV^2 = \text{constant} \). To find \( P \) as a function of \( V \), we can rewrite the equation as:
\[
P = \frac{P_{\text{initial}} \times V_{\text{initial}}^2}{V^2}
\]
Performing the integration, the integrand \(\frac{1}{V^2}\) integrates to \(-\frac{1}{V}\). Therefore, evaluating the integral from \( V_{\text{initial}} \) to \( V_{\text{final}} \) yields:
\[
W = P_{\text{initial}} \times V_{\text{initial}}^2 \left[ -\frac{1}{V} \right]_{V_{\text{initial}}}^{V_{\text{final}}}
\]
Substituting \( P_{\text{initial}} = 1 \, \text{bar} \) and \( V_{\text{initial}} = 0.1 \, \text{m}^3 \), we get:
\[
W = 1 \, \text{bar} \times 0.1 \, \text{m}^3 \left( -\frac{1}{0.0333} + \frac{1}{0.1} \right)
\]
Simplifying this yields:
\[
W = -2 \, \text{kJ}
\]

\newpage
\section*{Problem \#4}
Gaseous \(\text{CO}_2\) is contained in a vertical piston–cylinder assembly by a piston of mass \(51 \, \text{kg}\) and having a face area of \(0.01 \, \text{m}^2\). The mass of the \(\text{CO}_2\) is \(50 \, \text{g}\). The \(\text{CO}_2\) initially occupies a volume of \(0.005 \, \text{m}^3\) and has a change in specific internal energy of \(\Delta u = 67 \, \text{kJ/kg}\). The atmosphere exerts a pressure of \(95 \, \text{kPa}\) on the top of the piston. Heat transfer occurs slowly from the \(\text{CO}_2\) to the piston assembly, and the volume of the \(\text{CO}_2\) decreases to \(0.0025 \, \text{m}^3\). Friction between the piston and the cylinder wall can be neglected. The change in kinetic and potential energy of the gas is negligible. Kinetic energy can be ignored, but potential energy should be considered. Considering the piston and gas together as a system in the cylinder, calculate the following:

\begin{enumerate}
    \item What is the change in specific internal energy of \(\text{CO}_2\) gas?
    \item What is the change in internal energy of the piston in kJ?
    \item What is the change in potential energy of the piston in J?
    \item What is the work done on the gas in kJ?
    \item What is the amount of heat \(Q\) released from \(\text{CO}_2\) to surroundings in kJ?
\end{enumerate}

\section*{Solution}

\subsection*{Change in Specific Internal Energy of \(\text{CO}_2\) Gas}
The change in specific internal energy \(\Delta u\) is given in the problem as \(67 \, \text{kJ/kg}\). Since the problem is asking for the specific internal energy, there is no need to multiply this by the mass of the \(\text{CO}_2\). Thus,

\[
\Delta u = 67 \, \text{kJ/kg}
\]

\subsection*{Change in Internal Energy of the Piston}
The internal energy of the piston would change by the same amount that the \(\text{CO}_2\) gas does, considering that heat transfer occurs between the gas and the piston assembly. Given that the system consists of both the piston and the gas, the energy lost by the gas is effectively transferred to the piston. Therefore, the change in internal energy of the piston is also

\[
\Delta U_{\text{CO}_2} = \Delta U_{\text{Piston}} = 0.05 \, \text{kg} \times 67 \, \text{kJ/kg} = 3.35 \, \text{kJ}
\]


\subsection*{Change in Potential Energy of the Piston}
The change in height \(\Delta h\) of the piston can be found using the change in volume and the face area \(A\):

\[
\Delta h = \frac{\Delta V}{A} = \frac{0.005 \, \text{m}^3 - 0.0025 \, \text{m}^3}{0.01 \, \text{m}^2} = 0.25 \, \text{m}
\]
The change in potential energy \(\Delta PE\) can be calculated as:

\[
\Delta PE = m \times g \times \Delta h = 51 \, \text{kg} \times 9.77 \, \text{m/s}^2 \times 0.25 \, \text{m} = 124.48725 \, \text{J}
\]

\subsection*{Work Done on the Gas}
The work done on the gas \(W\) can be calculated using the formula:

\[
W = P \times \Delta V = 95 \, \text{kPa} \times (0.0025 \, \text{m}^3 - 0.005 \, \text{m}^3) = -0.2375 \, \text{kJ}
\]

\subsection*{Amount of Heat \(Q\) Released}
Using the first law of thermodynamics:

\[
Q = \Delta U_{\text{Total}} + W
\]
The total change in internal energy \(\Delta U_{\text{Total}}\) is the sum of the change in internal energy of the gas and the change in potential energy of the piston:

\[
\Delta U_{\text{Total}} = \Delta U_{\text{CO}_2} + \Delta PE = 3.35 \, \text{kJ} + 0.12448725 \, \text{kJ} = 3.47448725 \, \text{kJ}
\]
Therefore, \(Q\) can be calculated as:

\[
Q = 3.47448725 \, \text{kJ} - 0.2375 \, \text{kJ} = 3.23698725 \, \text{kJ}
\]
\newpage
\section*{Problem \#5}
A refrigeration cycle is operating as shown in the figure with a coefficient of performance \(\beta = 3.0\). For this cycle, \(Q_{\text{out}} = 300 \, \text{kJ}\). We are tasked to find:
\begin{enumerate}
    \item The heat input \(Q_{\text{in}}\) in kJ.
    \item The work of the cycle \(W_{\text{cycle}}\) in kJ.
\end{enumerate}

\section*{Solution}

\subsubsection*{Calculating \(W_{\text{cycle}}\)}
The correct formula for the coefficient of performance \(\beta\) in a refrigeration cycle, according to the problem, is:

\[
\beta = \frac{Q_{\text{out}} - W_{\text{cycle}}}{W_{\text{cycle}}}
\]

We are given that \(\beta = 3.0\) and \(Q_{\text{out}} = 300 \, \text{kJ}\). Solving for \(W_{\text{cycle}}\):

\[
W_{\text{cycle}} = \frac{Q_{\text{out}}}{\beta + 1} = \frac{300 \, \text{kJ}}{4} = 75 \, \text{kJ}
\]

\subsubsection*{Calculating \(Q_{\text{in}}\)}
To find \(Q_{\text{in}}\), we can use the energy balance equation:

\[
Q_{\text{in}} = Q_{\text{out}} - W_{\text{cycle}}
\]

\[
Q_{\text{in}} = 300 \, \text{kJ} - 75 \, \text{kJ} = 225 \, \text{kJ}
\]

\subsection*{Summary}
\begin{enumerate}
    \item The heat input \(Q_{\text{in}}\) for the refrigeration cycle is \(225 \, \text{kJ}\).
    \item The work of the cycle \(W_{\text{cycle}}\) is \(75 \, \text{kJ}\).
\end{enumerate}

\end{document}
